\section{결론}
\label{sec:conclusion}

본 논문은 연속 CoC에서 이전 사건의 요구, 요구가 끝나는 조건, 행동 순서, 그리고 참ㆍ거짓ㆍ정보 부족으로 나뉘는 증거를 읽어 검사 규칙으로 만들었다. 이 규칙은 이전 요구를 기억하는 Python 검사기와 같은 방식으로 동작하는 \uppaal 검사 모델에서 실행했다. 중심 결론은 분명하다. 본 연구에서 정한 규칙과 합성 시험 사례에서는 이전 요구를 기억해야 ``정지 요구가 끝나기 전에 진행한 경우''와 ``끝난 요구를 검사 목록에서 지우지 않은 경우''를 모두 구분할 수 있다.

가장 직접적인 근거는 정상 사례와 한 항목을 바꾼 사례로 이루어진 합성 시험 40쌍이다. 이전 요구를 기억하는 검사는 바꾼 사례 40/40에서 미리 지정한 오류를 판정했고, 현재 사건만 보는 검사는 20/40만 찾았다. 정상 사례를 오류로 잘못 판정한 경우는 두 검사 모두 0/40이었다. 두 검사의 차이인 20건은 모두 앞 문장의 요구를 알아야 판단할 수 있는 사례였다. 즉, 요구가 끝나기 전에 진행했거나 끝난 요구를 검사 목록에서 지우지 않은 경우였다. 또한 Python과 UPPAAL의 7개 표준 시험 사례에서는 최종 판정, 오류 종류, 최초 위반 사건과 정보 부족 사유가 모두 같아 불일치가 0/7이었다. 이는 두 구현이 시험한 범위에서 같은 판정과 위반 기록을 냈음을 보여 준다.

실제 주행 자료에서는 90개 장면에서 이어지는 사건 두 개씩을 묶은 309개 구간을 만들고, 그중 27개 장면을 검토했다. 검토자들의 첫 판단이 얼마나 비슷한지를 나타내는 Fleiss $\kappa$는 0.027로 낮았고, 서로 다른 판단의 이유를 논의한 뒤에는 27개 모두에서 모순이 없다는 데 합의했다. 이는 실제 오류의 비율을 계산한 결과가 아니라, 자동 생성 설명이 모호할 수 있으며 사람 사이의 합의 과정이 필요함을 보여 주는 결과이다.

보조 분석에서는 \texttt{scene-0001}에서 한 항목씩 바꾼 네 사례가 모두 예상한 검사 조건에서 실패했다. 행동 방향이 분명한 134개 사건은 기본 설정에서 속도 변화 확인 125건, 사건 시점에 이미 요구를 만족한 4건, 추가 검토가 필요한 5건으로 나뉘었다. 제한 시간을 $D=3$초로 두었을 때 두 장면은 첫 판단 시각부터 시간을 재면 실패하고 반복 판단 시각부터 다시 재면 통과했다. 속도로 반응을 판단한 130건 중 58건은 27개 속도 판정 설정에서 같은 결과를 유지하지 않았다. 이와 별도의 장면 경계 시험은 연결 근거 없이 앞 장면의 요구를 다음 장면으로 넘기는 오류를 구분했다. 이 결과는 검사 조건, 속도에서 계산한 증거, 장면 연결 근거를 함께 관리해야 함을 보여 주지만, 실제 자료의 오류를 얼마나 잘 찾는지를 직접 측정한 것은 아니다.

종합하면 본 결과는 여러 사건에 걸쳐 요구가 이어질 때 이전 요구를 기억하는 검사가 필요하고 효과적임을 뒷받침한다. 또한 같은 검사 규칙을 Python과 UPPAAL로 실행할 수 있으며, 시험 사례에서는 두 구현의 최종 판정과 위반 기록이 같음을 확인했다. 다음 단계는 별도로 수집한 실제 오류 자료에 객체 상태와 이동 경로 증거를 더하고 학습 전후 성능을 비교하여, 이 효과가 실제 오류 검출과 더 정확한 행동 학습으로 이어지는지를 평가하는 것이다.
